\documentclass[11pt, a4paper]{article}
\usepackage[utf8]{inputenc}
\usepackage{amsmath}
\usepackage{amssymb}
\usepackage{geometry}
\geometry{a4paper, margin=1in}
\usepackage{booktabs}

% === 新增的行距與段落間距設定 ===
\usepackage{setspace}
\onehalfspacing           % 設定 1.5 倍行距
\setlength{\parskip}{1em}
% 設定段落間距
% ================================



\title{Operations Research, Spring 2026 (114-2) \\ Final Project Proposal}

\author{
    \textbf{Group 6} \\[0.4cm] % 填入組別，並在下方增加 0.4cm 的間距
    R14725049 Yi-Chen, Tsai \\[0.4cm]
    R14323011 Pei-Hsuan, Chung \\[0.4cm]
    R13723065 Cheng-Hsien, Li \\[0.4cm]
    R13723045 Yung-Yao, Li	 \\[0.4cm]
    B09302150 Chien-Ting, Cheng\\[0.4cm]
    National Taiwan University \\[0.4cm]
}

\date{April 04, 2026}

\begin{document}

\maketitle

\section{Introduction}

In the operational framework of professional baseball, particularly in Major League Baseball (MLB) with its grueling 162-game schedule, the allocation and scheduling of pitching resources are undoubtedly the cornerstones of a team's success.
Pitchers represent a team's most valuable yet fragile assets, as their physiological limits and injury risks directly dictate performance over a high-intensity regular season.
To maintain competitiveness in such a demanding environment, a team must precisely manage appearance frequencies and pitch counts to ensure these scarce resources are deployed at optimal moments.
This complex challenge of multi-objective trade-offs and temporal resource distribution represents a core area where Operations Research (OR) provides significant value to sports management.

Typical pitcher scheduling involves highly intricate physiological and institutional constraints requiring differentiated management for various roles: starting pitchers responsible for stabilizing innings, middle relievers bridging the gap mid-game, and relief pitchers (closers) tasked with securing victories. In practice, teams must adhere to strict conventions; for instance, starters typically follow a "start on four days’ rest" or "five days’ rest" rotation to maintain arm health. Conversely, middle and relief pitchers must account for limits on consecutive appearance days and weekly innings to mitigate the risk of overuse injuries. While these conventions provide a basic protective framework, optimizing them within a global scheduling model—especially amidst MLB’s dense schedule and frequent cross-city travel—remains a formidable analytical challenge.

Generally, this decision-making process can be divided into two stages.
The first stage involves determining the official roster and player roles based on valuation and talent profiles, while the second stage focuses on crafting a detailed personnel schedule for the specific season itinerary.
Although coaching staffs must dynamically adjust their strategies during actual gameplay due to opposing lineups or other uncertainties, many organizations currently lack a scientific baseline tool.
Without a long-term strategic anchor, teams often find themselves in a reactive position, making isolated decisions that may compromise the team's health and performance later in the season.

This study aims to assist teams in optimizing the second-stage decision process by establishing an `initial default' pitcher scheduling baseline before the season begins and before in-season variables occur.
To simplify the problem while preserving the core modeling logic, this study assumes that the expected runs scored in each game are given, estimated in advance using external analytical tools.
Under this setting, the model focuses on optimizing pitcher deployment across the season to maximize the team's expected number of wins (or win rate), while strictly adhering to physiological constraints such as the “start on four days’ rest” rule.
This baseline schedule serves not only as a strategic roadmap for the season but also provides a quantitative coordinate for the team to rationally evaluate the long-term impact of tactical adjustments when unexpected system perturbations occur.

\section{Problem description}

To strategically assign a roster of $P$ pitchers throughout a full season of $G$ games and maximize the total number of wins, this model assumes that the coaching staff has pre-obtained simulated scoring data ($C_k$) for the team's lineup in each game. The primary objective is to utilize historical performance data (expected runs allowed per slot, $E_{ij}$) and incorporate pitcher protection protocols to determine the optimal assignment for each game.

The total runs allowed in a game are calculated by summing the expected runs from the pitchers assigned to each slot. A game is recorded as a win ($w_k = 1$) if the total runs allowed are less than the runs scored by the lineup; otherwise, it is a loss ($w_k = 0$).

In this model, a standard 9-inning game is divided into four specific time slots:
\begin{itemize}
    \item $j=1$: Starting Pitcher (Innings 1--5).
    \item $j=2$: Middle Relievers (Innings 6--7).
    \item $j=3, 4$: Closer (Innings 8 and 9).
\end{itemize}


\textbf{Note on Late-Inning Roles:} In standard MLB operations, the 8th inning is typically handled by a Setup Man, while the 9th inning is reserved for the Closer. However, because both roles operate under high-leverage situations and share nearly identical physiological demands, usage patterns, and rest requirements, our model uniformly classifies pitchers assigned to the 8th ($j=3$) and 9th ($j=4$) innings under the ``Closer'' category. Consequently, they are subject to the exact same set of protective constraints.

Each pitcher is restricted to at most one time slot per game.
Furthermore, the following pitcher protection constraints are enforced:
\begin{enumerate}
    \item \textbf{Starting Pitcher Rest}: A pitcher assigned as a starter ($j=1$) cannot appear in the subsequent four games.
    \item \textbf{Closer Fatigue}: A pitcher assigned as a closer ($j=3, 4$) cannot pitch in three consecutive days.
    \item \textbf{Reliever Rest}: Middle relievers ($j=2$) are prohibited from appearing in the immediate next game.
    \item \textbf{Seasonal Limit}: No pitcher can exceed 25 total starts in a single season.
\end{enumerate}

\begin{table}[h]
\centering
\caption{Mathematical Notations for the Pitcher Scheduling Model}
\begin{tabular}{ll}
\hline
\textbf{Symbol} & \textbf{Description} \\ \hline
\textit{\textbf{Sets and Indices}} & \\
$I = \{1, \dots, P\}$ & Set of pitchers, indexed by $i$. \\
$J = \{1, 2, 3, 4\}$ & Set of roles/slots: 1 (Starter), 2 (Relievers), 3-4 (Closer). \\
$K = \{1, \dots, G\}$ & Set of games in the season, indexed by $k$. \\ \hline
\textit{\textbf{Parameters}} & \\
$G$ & Total number of games (e.g., 162). \\
$P$ & Total number of pitchers (e.g., 12). \\
$C_k$ & Expected runs scored by the lineup in game $k$. \\
$E_{ij}$ & Expected runs allowed by pitcher $i$ in slot $j$. \\
$M$ & A large positive constant for Big-M logic. \\ \hline
\textit{\textbf{Decision Variables}} & \\
$x_{ijk} \in \{0, 1\}$ & 1 if pitcher $i$ is assigned to slot $j$ in game $k$; 0 otherwise. \\
$w_k \in \{0, 1\}$ & 1 if game $k$ is a win; 0 otherwise. \\ \hline
\end{tabular}
\end{table}

\section{Mathematical model}

\subsection*{Objective Function}
With these notations, the total wins that the team seeks to maximize is:
\begin{equation}
\max \sum_{k\in K}w_{k}
\end{equation}

\subsection*{Constraints}
The following constraints ensure valid pitcher assignments and adhere to protection protocols:

\begin{align}
& \sum_{i\in I}\sum_{j\in J}E_{ij}x_{ijk} - C_{k} \le M(1-w_{k}) && \forall k\in K \\
& \sum_{j\in J}x_{ijk} \le 1 && \forall i\in I, k\in K \\
& \sum_{i\in I}x_{ijk} = 1 && \forall j\in J, k\in K \\
& \sum_{t=0}^{2}\sum_{j\in J}x_{i,j,k+t} \le 2 && \forall i\in I, k\in K\backslash\{G-1,G\} \\
& x_{i2k} + \sum_{j\in J}x_{ij,k+1} \le 1 && \forall i\in I, k\in K\backslash\{G\} \\
& x_{i1k} + \sum_{t=1}^{4}\sum_{j\in J}x_{ij,k+t} \le 1 && \forall i\in I, k\in K\backslash\{G-3,...,G\} \\
& \sum_{k\in K}x_{i1k} \le 25 && \forall i\in I
\end{align}

The feasible region of the proposed model is governed by several sets of critical constraints.
Specifically, Constraint (2) utilizes the Big-M method to link the binary winning variable $w_{k}$ with the scoring performance, ensuring that $w_{k}$ is assigned a value of 1 only if the total expected runs allowed by the assigned pitchers do not exceed the team's offensive runs $C_{k}$.

Regarding the assignment logic, Constraint (3) mandates that each pitcher is assigned to at most one time slot in any single game to prevent overlapping roles, while Constraint (4) ensures that every designated time slot is filled by exactly one pitcher to satisfy tactical requirements for a complete game.

To manage pitcher health and workload, Constraint (5) prohibits closers from appearing in more than two games within any three-day window, effectively preventing three consecutive days of pitching.
While the restriction can be universally applied to all pitcher roles including starting pitchers and middle relievers as the rest protocols for those pitchers are significantly more stringent.

Furthermore, Constraint (6) enforces a recovery period for middle relievers $(j=2)$ by prohibiting immediate next-game assignments.
For starting pitchers $(j=1)$, Constraint (7) mandates a minimum four-game rest period in accordance with professional health protocols.
Finally, Constraint (8) imposes a seasonal upper limit of 25 starts per pitcher to ensure long-term physical sustainability and workload management.

\textbf{Calibration of Big-M ($M$):}

To properly enforce the logical implication in Constraint (2) using the Big-M method, the parameter $M$ must serve as a valid upper bound for the left-hand side of the inequality: $\sum_{i\in I}\sum_{j\in J} E_{ij}x_{ijk} - C_{k}$.

A mathematically rigorous upper bound must evaluate the extreme values of both terms independently.
Since the decision variables $x_{ijk}$ are binary, the supremum of the first term occurs when all pitchers are theoretically assigned simultaneously, yielding $\sum_{i,j} E_{ij}$.
For the second term, because a baseball team cannot score negative runs, its minimum possible value is strictly bounded by $C^{min}_k = 0$.
Consequently, a theoretically safe upper bound for the entire expression is derived as $\sum_{i,j} E_{ij} - 0$.

However, governed by Constraint (4), exactly four pitcher slots are filled per game, meaning the true mathematical upper bound is significantly lower than this total sum.
Furthermore, employing such a loose theoretical bound often results in weak linear relaxations, which can substantially deteriorate the solver's computational efficiency.
To establish a tighter, practical bound, we integrate empirical domain knowledge.
The modern MLB record for the largest single-game run differential occurred in 2007 (Texas Rangers vs. Baltimore Orioles, 30--3), establishing an extreme historical margin of 27 runs.
Grounded in this worst-case scenario, setting $M$ to a constant such as 40 or 50 provides a sufficient buffer to encompass any realistic baseball anomaly while preserving a tight feasible region for optimal solver performance.

\section{Expected results}
Our proposed model will evaluate pitcher performance using Earned Run Average (ERA), leveraging historical data sourced from leading databases such as Baseball Savant, FanGraphs, and Baseball Reference.

The primary objective is to optimize pitcher allocation across a season, subject to the realistic operational constraints detailed in Section 2, in order to maximize a team's expected wins.
Given the computational complexity of scheduling for a standard 162-game season, relying solely on exact Integer Programming (IP) will be intractable.
Therefore, we aim to implement heuristic algorithms designed to efficiently generate near-optimal schedules.
To validate our approach, we will benchmark our heuristic results against exact IP solutions on smaller problem instances.

We anticipate achieving a minimal optimality gap and outperforming existing baseline heuristics.
Ultimately, we envision this research yielding a robust, computationally efficient framework that translates complex mathematical optimization into highly actionable, real-world strategies for front offices.
\end{document}
